\chapter{Intuitive Derivations}
\label{app:intuition}

\section{Bells as Waves}
\label{app:waves}

The Hubbard-Stratonovich identity writes a bell as an expectation of $\exp(-iu(z - c))$ over $u \sim \mathcal{N}(0, \gamma)$, with $c = \frac{1}{2\gamma}$.
To see the waves, split each term of the expectation with Euler's formula $e^{-ix} = \cos x - i \sin x$:
\[
  \exp \left(-iu \left(z - c\right) \right) = \cos \big(u (z - c)\big) - i \sin \big(u (z - c)\big).
\]
The sine is odd in $u$, and since $u$ and $-u$ are equally likely under $\mathcal{N}(0, \gamma)$, it cancels and only the cosine remains:
\[
  \exp \left(\frac{1}{2}z - \frac{1}{2}\gamma z^2 \right) = \exp \left(\frac{1}{8\gamma}\right) \E_{u \sim \mathcal{N}(0, \gamma)} \left[ \cos \big(u (z - c)\big) \right].
\]

\section{The Characteristic Function as an Area}
\label{app:cf-area}

The characteristic function is used to compute $\E_{p(\xm \mid \xo)} \left[ \exp \left(-iu \left(z - c\right) \right) \right]$.
The integrand depends on $\xm$ only through $z = w^T x$, so we can take the expectation over the distribution $p(z)$ of $z$ instead.
The real part of $\exp(-iu(z - c))$ is the wave $\cos(u(z - c))$ from Appendix~\ref{app:waves}:
\[
  \mathrm{Re} \left[ \E_{p(\xm \mid \xo)} \left[ \exp \left(-iu \left(z - c\right) \right) \right] \right] = \E_{p(z)} \left[ \cos \big(u (z - c)\big) \right] = \int \cos \big(u (z - c)\big)\, p(z) \, \mathrm{d}z.
\]
This is the area under the wave times $p(z)$, where parts below zero count negatively.
So for every frequency $u$, the characteristic function gives the signed area under the wave times $p(z)$.

\section{Other Rules for the PG Integral}
\label{app:pg-rules}

Besides Gauss--PG (Section~\ref{subsec:pg-quadrature}), we consider three
families of rules for the outer integral over $\gamma \sim \PG(1,0)$: trapezoid,
Gauss--Legendre and Gauss quadrature, each with the nodes spaced in $\gamma$ and in
$\log \gamma$, and rules whose weights are fitted to the sigmoid. All rules
except Gauss--PG need the density $p$ of $\PG(1,0)$ pointwise. It has no closed
form, but two alternating series, one converging fast for small and one for large
$\gamma$, evaluate it to machine precision. All weights are normalised to sum to
one.

\subsection{Spaced in $\gamma$ and in $\log \gamma$}

\paragraph{Trapezoid.}
The composite trapezoid rule on $[0, T]$ with $n$ equally spaced nodes $\gamma_j$
and spacing $h$ has the weights $w_j \propto c_j\, p(\gamma_j)\, h$, with
$c_j = \frac{1}{2}$ at both ends and $1$ otherwise. We use $T = 5$, above which
$\PG(1,0)$ has a mass of only about $2 \cdot 10^{-11}$. With equal spacing, only
few nodes fall where the mass is, around $\gamma \approx 0.1$ to $0.5$.

\paragraph{Log-trapezoid.}
Substituting $\gamma = e^s$ turns the expectation into
\[
  \E_{\gamma \sim \PG(1,0)}[f(\gamma)] = \int_{-\infty}^{\infty} f(e^s)\, p(e^s)\, e^s \,\mathrm{d}s.
\]
Applying the trapezoid rule with $n$ equally spaced nodes $s_j$ and spacing $h$
gives nodes $\gamma_j = e^{s_j}$ and weights $w_j \propto p(\gamma_j)\,\gamma_j\,h$.
Since $p$ decays like $e^{-1/(8\gamma)}$ towards $0$ and like $e^{-\pi^2\gamma/2}$
towards infinity, the range $\gamma \in [10^{-3}, 10]$ covers all of the mass.
Equal spacing in $\log\gamma$ places many nodes at small $\gamma$, exactly where
$e^{-\gamma z^2/2}$ lives for large $|z|$, and for smooth integrands that vanish at
both ends the trapezoid rule converges exponentially in $n$.

For the sigmoid alone, this makes the log-trapezoid rule more accurate than
Gauss--PG with the same number of nodes. Inside PCPG, however, every node comes
with the factor $e^{1/(8\gamma_j)}$ of $f$, which multiplies the error of the
inner expectation. At the smallest node $\gamma = 10^{-3}$ this factor is about
$10^{54}$, which overflows in floating point and makes the estimate useless. The
smallest Gauss--PG node for $n = 19$ is $\gamma \approx 0.046$, where the factor
is below $16$.

\paragraph{Gauss--Legendre.}
The standard Gauss--Legendre nodes $x_j \in [-1, 1]$ with weights $v_j$ are mapped
to $\gamma_j = T (x_j + 1)/2$ on the same interval $[0, T]$, and the density goes
into the weights, $w_j \propto v_j\, \frac{T}{2}\, p(\gamma_j)$. The rule is exact
for polynomials times the density, which is not what we integrate, so it does not
reach the convergence rate of a Gaussian rule for $\E[f(\gamma)]$.

\paragraph{Log-Gauss--Legendre.}
The same in $s = \log \gamma$, on the same range as the log-trapezoid rule: the
Legendre nodes are mapped to $[\log 10^{-3}, \log 10]$, and the weights get the
factor $p(e^s)\, e^s$ as there.

\paragraph{Log-Gauss--PG.}
The Gaussian quadrature rule in $s = \log \gamma$, with the distribution of $s$ as
the weight function, so it is exact for polynomials in $s$ up to degree $2n - 1$.
The moments of $s$ have no closed form, so the Jacobi matrix is not built from
moments but by the Lanczos (Stieltjes) procedure on a fine discretisation of the
distribution of $s$, with 4000 trapezoid points on $[\log 10^{-4}, \log 20]$.
Lanczos builds the orthonormal polynomials point by point, with
reorthogonalisation so that it stays stable, and their recurrence coefficients
form the Jacobi matrix. The nodes and weights then follow from its
eigendecomposition as in the Golub--Welsch algorithm.

\subsection{Fitting the Nodes to the Sigmoid}
\label{subsubsec:sigmoid-fit}

With the sign flip (Section~\ref{subsec:sign-flip}), the rule is only evaluated
for $z$ around and below $0$, and the error that matters is
$\tfrac{1}{2}e^{z/2}\varepsilon(z^2/2)$ on that side. None of the rules above
optimise this quantity, since they all discretise the $\PG(1,0)$ distribution.

Since the rule enters the estimator linearly, the estimator computes exactly
$\E[\hat\sigma(z)]$, and its error is the error of $\hat\sigma$ averaged over the
distribution of $z$. Nothing in the rest of the method requires the nodes to
represent $\PG(1,0)$: HS and the characteristic function only need $\gamma_j > 0$,
and the weights only need to be nonnegative. We can therefore choose the nodes and
weights to approximate the sigmoid itself. On a grid of 400 log-spaced candidate
nodes $\gamma_j \in [0.02, 8]$ and 1500 fitting points $z_r \in [-30, 5]$, we
minimise a norm of the error curve
\[
  e_r = \tfrac{1}{2}\, e^{z_r/2} \sum_{j} w_j\, e^{-\gamma_j z_r^2/2} - \sigma(z_r)
\]
over nonnegative weights that sum to one. We use three norms:
\begin{itemize}
  \item \textbf{L2}, the sum of squared errors, solved by nonnegative least squares with the sum-to-one constraint as an extra row,
  \item \textbf{L1}, the sum of absolute errors, i.e.\ the area under the error curve,
  \item \textbf{max}, the largest absolute error, which makes the error curve equally high everywhere.
\end{itemize}
L1 and max are linear programs. All three keep only a few of the candidates, so the
number of nodes is not chosen directly. The sum-to-one constraint makes the
estimate exactly $\frac{1}{2}$ at $z = 0$. The margin up to $z = 5$ covers the part
of the distribution of $z$ that remains on the positive side after the flip, and
the grid starts at $\gamma = 0.02$, so that $e^{1/(8\gamma)}$ in the inner integral
stays below about $500$.

This changes the view of the outer integral: the method approximates the
sigmoid by functions of the form $\tfrac{1}{2}e^{z/2}e^{-\gamma z^2/2}$, whose
expectations are tractable, and the PG identity is one particular, exact but
infinite, choice of how to combine them. A finite rule only has to combine them
well, not to represent the PG distribution. The only variants this does not apply
to are the ones that sample $\gamma$ from $\PG(1,0)$, since a fitted rule cannot
be sampled from.

\section{Variants of PCPG}
\label{app:variants}

Each of the two expectations of PCPG can be computed by quadrature or by Monte
Carlo, which gives the four combinations in Table~\ref{tab:pcpg-variants}. The
variants with Monte Carlo use Algorithm~\ref{alg:pcpg} with the sampled nodes
and the equal weights $\frac{1}{M}$ and $\frac{1}{N}$ in place of the quadrature
nodes and weights, and without the factor $\frac{1}{\sqrt{\pi}}$ of
Gauss--Hermite.

\begin{table}[H]
\centering
\small
\begin{tabular}{lll}
\hline
\textbf{Variant} & \textbf{Outer} $(\gamma_j, w_j)$ & \textbf{Inner} $(u_{ij}, \omega_{ij})$ \\
\hline
PCPG-MC   & $\gamma_j \sim \PG(1,0)$, $w_j = \frac{1}{M}$ & $u_{ij} \sim \mathcal{N}(0,\gamma_j)$, $\omega_{ij} = \frac{1}{N}$ \\[4pt]
PCPG-GH   & $\gamma_j \sim \PG(1,0)$, $w_j = \frac{1}{M}$ & Gauss--Hermite: $u_{ij} = \sqrt{2\gamma_j}\,t_i$, $\omega_{ij} = \frac{\omega_i}{\sqrt{\pi}}$ \\[4pt]
PCPG-GPG  & Gauss--PG rule & $u_{ij} \sim \mathcal{N}(0,\gamma_j)$, $\omega_{ij} = \frac{1}{N}$ \\[4pt]
PCPG      & Gauss--PG rule & Gauss--Hermite: $u_{ij} = \sqrt{2\gamma_j}\,t_i$, $\omega_{ij} = \frac{\omega_i}{\sqrt{\pi}}$ \\
\hline
\end{tabular}
\caption{The four combinations of Monte Carlo and quadrature for the two expectations of PCPG. Hermite nodes $(t_i, \omega_i)$ as in Section~\ref{subsec:gauss-hermite}, Gauss--PG nodes $(\gamma_j, w_j)$ as in Section~\ref{subsec:pg-quadrature}.}
\label{tab:pcpg-variants}
\end{table}

PCPG-GPG samples the inner expectation, whose variance is very high, so it
performs much worse than PCPG in practice.

\section{Why Large Positive $z$ Is Hard}
\label{app:large-z}

In \eqref{eq:sigma-laplace}, any error $\varepsilon$ a rule makes in approximating
$\mathcal{L}(z^2/2)$ enters the estimate of $\sigma(z)$ multiplied by
$\tfrac{1}{2}e^{z/2}$. The nodes only see $z^2$, so this error is the same for $z$
and $-z$, but its effect is not:
\[
  \bigl|\hat\sigma(z) - \sigma(z)\bigr| = \tfrac{1}{2}\, e^{z/2}\, \varepsilon(z^2/2).
\]
For $z < 0$ the factor shrinks the error, for $z > 0$ it amplifies it
exponentially. Staying at $\sigma(z) \approx 1$ for large positive $z$ requires the
sum of exponentials to reproduce the tail $\mathcal{L}(z^2/2) \approx 2e^{-z/2}$
almost exactly, which only rules with nodes very close to $\gamma = 0$ can do.
Rules with few or no nodes there (trapezoid, Gauss--Legendre, Gauss--PG) have a
tail that is too small, so $\hat\sigma(z)$ falls back to $0$ and the error approaches
$1$. A rule that puts more weight at tiny $\gamma$ than $\PG(1,0)$ really has there
does the opposite: its tail is too heavy, and $\hat\sigma(z)$ blows up like $e^{z/2}$.
Only the rules spaced in $\log \gamma$ stay accurate without further help, and only
with enough nodes.
